\section{Discussion and Limitations}
\label{sec:discussion}

Under the largest-first rule the four metrics agree on when each site is weakest: all reach
their minimum in the minimum-commitment hours, and no metric tells more than the committed set
of synchronous generators. With every unit online they disagree, although each then varies by
only a few percent over the week. The two rules, however,
place the weakest hour at opposite ends of the load range. The critical hour is thus a property
of the commitment, not of the metric, and a screen over time needs a model of commitment or
measured operating states. A minimum-load snapshot finds the weekly minimum where, as under
the largest-first rule here, the fewest units are online at minimum load.

Under outages, commitment still sets the weakest hour in nearly every case, and the outage
changes how weak the site is rather than when. The metrics differ in how they report that
change. SCR and WSCR follow the change in the network.
ESCR and SDSCR add the other plants' contribution, weighted by their electrical proximity to the site,
and an outage that moves another plant away lowers that weight. All the points in which they
oppose the network come from two outages, each of a line at another plant's POI, and the
interaction term alone explains them. A screen based on ESCR or SDSCR would therefore not flag
these outages, while one based on SCR would. Whether the site is in fact weaker or stronger 
cannot be decided from fault quantities, since both readings follow from the same impedances. 
It depends on whether the site's stability is limited more by the network's fault current or 
by the interaction between nearby plants, which only an EMT study can show.

The opposition found above does not come from a convention but from the interaction term in the
metrics' definitions. How the interaction factor is obtained, however, changes how often ESCR
opposes the network and by how much. Since the literature obtains it in different ways \cite{cigre364,ieee2800,fingrid2023}, an ESCR value should state how its
interaction factor is obtained.

The findings rest on the following modelling choices. The fault model assumes constant-impedance loads and
constant-current converters, which makes the fault network linear and \eqref{eq:shape} close to
exact. Converters that inject reactive current during the fault would make it nonlinear and add
their own contribution to the short-circuit power. The constant-impedance loads are also the only
path by which $|Z_{ii}|$ changes with load when every unit is online, so that result depends on
the load model. The voltage-step factor uses ideal voltage-controlled generators, so it remains
untested whether its smaller interaction comes from the network or from the ideal control, and
its dependence on the step size is not explained. The study covers one network and one week,
with the IBR plants at rated output and unity power factor, switched shunts fixed, commitment
held during outages and single branch outages only. Pre-fault voltages were not corrected to an
operating band. Since every plant runs at rated output, denominators based on rated, actual or
apparent power differ only by one constant factor per site and change no conclusion. The 1\,\%
threshold for a material change is a choice. The metric equations follow their sources, but the
computed values could not be checked against a published case.
